% Section 8. Drafted from docs/PAPER-OUTLINE.md section 8 and docs/THEORY.md "Interfaces between the
% layers". Every named result is a gated, axiom-free theorem in coq/*.v (LayerInterfaces.v unless noted).

\section{The interfaces}\label{sec:interfaces}

The natural objection to a three-layer framework is that executions (A), consistent states (B) and
long-run behavior (C) are three classical theories placed side by side. This section answers it with
the three interfaces between the layers, each proved without changing the layers, and each with the
counterexample that marks where it stops.

% ---------------------------------------------------------------------------------------------
\subsection{A to C: Newman as the bridge, and where fairness falls short}

In layer~C's terms a reduction is a step of some action that changes the state, and its normal forms
are the settled states (\coq{nf_iff_stable}). Under termination from the start and decidable
normality, local joinability at reachable peaks, confluence and unique normal forms coincide
(\coq{newman_bridge}); the converses need no termination (\coq{cr_peaks}, \coq{cr_unf}). This is
Newman's lemma~\cite{newman1942} restricted to the reachable states. The single-registry theorem is
its instance, with termination of repair supplying the hypothesis and CC2 being state descent.

Fair settlement does not replace termination. In the standard four-point system
$k_a \leftarrow k_b \leftrightarrow k_c \rightarrow k_d$, every peak joins and every fair run settles,
yet the run that only swaps $k_b$ and $k_c$ is an infinite reduction that starves the other actions, and
$k_a$ and $k_d$ are two normal forms below $k_b$ (\coq{four_point_fair_not_confluent}).

\begin{theorem}[\coq{fair_starves}, \coq{fair_sn_exact}]\label{thm:fair-sn}
If every fair run settles (C2), then every reachable loop of moving steps starves some action: some
action it never picks would move every state on the loop. Consequently, on a finite state space,
termination holds exactly when no reachable loop starves an action, and then Newman's lemma applies
(\coq{fair_newman}).
\end{theorem}

Fair settlement replaces termination exactly when reductions cannot be hidden forever by starving an
action. The condition is exact for termination, not for confluence: a network with a starved loop
can still be confluent (\coq{drop_starved_confluent}). Whether \cref{thm:fair-sn} is stated in the
fair-termination literature~\cite{lehmann1981,francez1986} is an open question of the review; the
statement is close to the standard picture, in which an infinite run is either fair or starves some
action forever.

% ---------------------------------------------------------------------------------------------
\subsection{B to C: the fixed points are the solutions}

On a store with the usual read and write laws, a state is fixed by every edge repair
$s[v := f(s_u)]$ if and only if it is a section, if and only if it solves the network's constraint
problem (\coq{edge_fixed_iff_section}, \coq{edge_fixed_iff_csp}); the same holds for resolver
repairs that recompute a vertex from several inputs (\coq{rupd_fixed_iff_csp}). Read in layer~C:

\begin{theorem}[\coq{singleton_iff_solution}, \coq{settled_canon_exact}]\label{thm:bc}
The fair recurrent classes that are single states are exactly the reachable solutions of the
constraint problem. If every fair run settles (C2), then every fair run settles on the result $c$ if
and only if every reachable solution lies in the fiber of $c$.
\end{theorem}

So a state that settles on the wrong result is a reachable solution outside the target fiber
(\coq{ghost_refutes}), and two reachable solutions with different results refute every target
(\coq{multistable}); a network with no solution has no raw settlement
(\coq{no_solution_no_settlement}). The converse direction fails in the way that matters most:

\begin{proposition}[\coq{flip2_unique_solution_livelock}]\label{prop:flip2}
There is a one-vertex network whose constraint problem has exactly one solution and whose repair
dynamics has a fair run that never settles.
\end{proposition}

A unique consistent state does not guarantee that runs reach it. This is the distinction the first
preregistered prediction (\cref{sec:tests-predictions}) bet on. One leak is recorded: a network with
no solution can still settle canonically, through a semantic oscillation whose repaired result never
changes (\coq{no_solution_canon_settles}).

% ---------------------------------------------------------------------------------------------
\subsection{A to B: consistency needs to hold only after repair}

A per-vertex family of maps that commutes with every edge map sends sections to sections
(\coq{natural_maps_sections}). For one network the exact condition is commutation on the values
sections take (\coq{preserves_exact}), and the commuting squares are necessary edge by edge for edges
that are not self-loops (\coq{square_necessary}), with counterexamples to the naive converse
(\coq{natural_converse_fails}). A uniform family is a unary polymorphism
(\coq{uniform_natural_is_unary_pol}).

For events in a network of registries, strict commutation with the repair is too strong. The
condition that matters is descent: an event must respect consistency only after repair.

\begin{theorem}[\coq{xu_is_descends}, \coq{candidate_implies_xu}, \coq{xu_iff_respects}]\label{thm:ab}
The federation's event condition XU is state descent along each registry's repair on its valid
states. Strict commutation implies it, and it holds exactly when each event passes to the quotient of
the valid states by the repair.
\end{theorem}

Strict commutation is sufficient but not necessary. On a supply federation XU holds, strict
commutation fails, and every order of events converges (\coq{strict_not_necessary}).
